\section{Further research questions}\label{sec:questions}

The leading constant is only the first part of the rounding problem.
The following questions separate plausible next steps from conclusions
already proved. They are proposed research problems; except for the
historical linear question, their open status refers to this report
and the sources inspected, not to an exhaustive classification of the
literature.

\begin{question}[Quantitative threshold errors]\label{q:rates}
For fixed $m>0$, find a usable rate in
\[
 T_m(K)=L_mK^{m+1}+E_m(K).
\]
Which exponents $\delta_m>0$ admit
$E_m(K)=O_m(K^{m+1-\delta_m})$? Can an explicit bound be obtained
uniformly for $m$ in a compact subinterval of $(0,\infty)$?
\end{question}

The proof identifies the correct place to work: make the crossing
times of the discrete quotient increments quantitative, then balance
their accumulated errors against the omitted small-time tail. The
classical $m=1$ estimate in \cref{eq:classical-error} supplies a
benchmark. Compactness alone contains no numerical rate.

\begin{question}[Bounded stopping-time error]\label{q:bounded}
For which $m>0$, if any, does
\[
 \tau_m(n)=(c_mn)^{1/(m+1)}+O_m(1)
\]
hold? In particular, resolve the stronger bounded-error claim for
$m=1$ recorded in A073047.
\end{question}

For pure powers, this is equivalent to
$T_m(K)=L_mK^{m+1}+O_m(K^m)$. One direction follows by inverting the
two neighboring threshold bounds. For the other, evaluate the
stopping-time formula at the first integer of each survival interval,
where the stopping index is $K+1$, and expand the resulting power.
The two equivalent formulations indicate that an error term at the
scale of a single threshold gap is required.

\begin{question}[Signs and fluctuation scales]
Does $T_m(K)-L_mK^{m+1}$ have an eventual sign? What are the natural
scales of its upper and lower envelopes? Does any centered and scaled
version have a limiting distribution along integer horizons?
\end{question}

The selected tabulated horizons have positive errors, while the denser
plotted data include both signs. Neither observation establishes an
eventual-sign theorem. The nonmonotone normalized errors show that a
simple monotonicity argument would need additional structure. Any probabilistic terminology here would describe a
distribution over sampled horizons, not randomness in the recurrence.

\begin{question}[A theory of microscopic schedules]\label{q:schedules}
Classify the limits for regularly varying weights when
$k(w_k/w_{k-1}-1)$ oscillates, has repeated zero values, or follows a
periodic pattern. Which information beyond the regular-variation
index determines the extinction constant and the limiting profile?
\end{question}

The block construction gives a completely solved test family. A
possible next class consists of a fixed periodic pattern of small
relative increments. The ceiling acts before any averaging, so an
effective drift would have to be derived from the individual phases
of that pattern. Establishing such an averaging theorem is a
substantive extension; it is not a consequence of
\cref{thm:weights-main}.

\begin{question}[Stability under sparse exceptional steps]
How many exceptional ratios can be allowed without changing $L_m$?
For example, suppose the local ratio condition holds outside a set
of indices of density zero. Which bounds on the exceptional
increments suffice, and can a sparse set of very large jumps change
the constant?
\end{question}

Finite modifications are already covered. Passing from finitely many
exceptions to a sparse infinite set requires controlling their
combined effect on the quotient increments. Density alone need not
control the size of those effects.

\begin{question}[Moving exponents and singular regimes]
Describe joint limits with $m=m(n)\downarrow0$ or $m=m(n)\to\infty$.
Determine ranges in which the fixed-parameter equivalent remains
uniform and identify any transition scales.
\end{question}

The discontinuity at $m=0$ makes the first question particularly
interesting: positive integer values never vanish when every weight
is exactly $1$. The explicit expansions in
\cref{prop:parameter} describe the constants, but they do not control
the speed of convergence of the underlying discrete process.

\begin{question}[Rates for the full trajectory and loss distribution]
Obtain explicit bounds for
\[
 \sup_{0\leq t\leq1}|X_n(t)-F_m(t)|
\]
and for convergence of the loss measures in
\cref{cor:losses}. Can one derive asymptotics for individual moments
or quantiles of the loss-time distribution with a controlled error?
\end{question}

The limiting density is explicit on each increment regime. It permits
certified computation of limiting observables using finite sums and
an integrable tail. Rates for the finite process again depend on
discrete crossing errors rather than solely on the Gamma constant.

\begin{question}[Sieve and game interpretations for higher powers]
Is there a natural sieve or solitaire model whose survival thresholds
are $T_m(K)$ for integer $m\geq2$? Can such a model explain the
Gamma product combinatorially or connect the threshold gaps to an
existing family of integer sequences?
\end{question}

The ordinary case has the Tchoukaillon interpretation. A higher-power
model could expose monotonicity, congruence, or conservation
properties not apparent in the floor recurrence. Any proposed model
should first be proved equivalent to the exact backward recurrence,
including its initial conditions.

\begin{question}[Faster exact algorithms and formal verification]
Can long consecutive ranges with a constant quotient increment be
skipped in a provably correct exact algorithm? Can the general
real-parameter theorem be formalized with a reusable library lemma
for monotone discontinuous drift limits?
\end{question}

The present algorithm takes one arithmetic step per stage. The
limiting piecewise linear structure suggests acceleration, but the
exact locations of discrete transitions must be certified before
steps can be skipped. For formalization, a useful first milestone is
the exact survival duality and rational-exponent algorithm, followed
by the compactness and endpoint-selection lemmas.

\subsection{A practical order for continuing the work}

The most direct continuation is a quantitative crossing-time analysis
for the pure power family. In parallel, a periodic-ratio model can
test how the constant changes when the local smoothness condition is
relaxed. These two directions address the principal limitations of
the present theorem: the absence of a finite error bound, and the
dependence on the microscopic schedule.

\appendix
\section{Proof dependencies and source audit}\label{app:audit}

\subsection{Logical dependencies}

The extinction proof proceeds through five independently checkable
steps:
\begin{enumerate}
\item Exact stagewise survival duality gives the backward ceiling
recurrence and the inverse threshold relation.
\item The integer increment bound and weighted rounding-error bound
give compactness and force positive terminal drift.
\item Strict monotonicity of $my(t)/t$ controls the ceiling
discontinuities, yielding a rigorous limit equation.
\item Solving its increment regimes gives the finite product;
the beta integral evaluates its asymptotic constant.
\item A telescoping estimate controls the entire omitted initial
part, allowing the two limits to be taken in the correct order.
\end{enumerate}
The smooth-weight theorem reuses these steps after proving the local
ratio and tail-sum consequences of its hypothesis. The forward-path
theorem follows from a second use of the exact stagewise duality.
No numerical estimate is an input to any of these proofs.

Two independent mathematical reviews checked the ceiling passage,
initial-value convention, smooth-weight tail, inversion, profile
endpoint argument, and block counterexample. The exact finite
verification was generated by a separate implementation. These
reviews improve confidence but are not a substitute for peer review
or formal verification.

\subsection{Sources reviewed and scope of the novelty statement}

The principal source records are the OEIS definitions and conjectures
for A073047, A082527, and A082528; the classical threshold entry
A002491; Erd\H{o}s and Jabotinsky's original sieve paper; the
Broline--Loeb preprint and the criticism attributed to Knuth in the
threshold entry; Brown's exposition; Leh\'ericy's 2023 proof for a
different generalized rounding rule; and Li's 2026 Tchoukaillon
preprint. The bibliography provides direct links.

Targeted searches combined the exact sequence identifiers with
power-rounding, asymptotic, Gamma, generalized sieve, and the
classical authors' names. Searches within ProveIt used the exact
sequence identifiers and Tchoukaillon-related terms. No matching
real-power theorem or formula for $c_m$ was found in that review.
This statement is intentionally bounded by the sources and search
scope. A conjectural label remaining on an OEIS entry does not prove
that no proof exists elsewhere.

The result should therefore be described precisely: this report
supplies a complete proof of the real-parameter conjecture as
recorded in A082528, identifies the square and cube constants, and
proves the stated extensions. It makes no claim to have newly
discovered the linear $\pi$ law or the general increment-regime
method. The numerical and symbolic evidence in the archive is
reproducible and separately labeled.

\section{Suggested OEIS update text}\label{app:oeis}

The following concise mathematical statements are suitable for
review before any submission. They have not been posted to OEIS.

\paragraph{For A082527.}
The conjectured asymptotic constant is
\[
c=2\Gamma(2/3)^3=4.965917162430455473\ldots.
\]
Thus $a(n)\asym(2\Gamma(2/3)^3n)^{1/3}$.
This is the $m=2$ case of the power-rounding extinction theorem.

\paragraph{For A082528.}
For every fixed real $m>0$, if $x_1=n$ and
$x_k=k^m\lfloor x_{k-1}/k^m\rfloor$, the first zero satisfies
\[
 a(n,m)\asym
 \left(m\Gamma\!\left(\frac{m}{m+1}\right)^{m+1}n\right)^{1/(m+1)}.
\]
In particular, the cubic constant is
$3\Gamma(3/4)^4=6.764822981008760234\ldots$.
A proof is given by the backward survival thresholds and their
piecewise linear scaling limit.

\paragraph{For the linear cross-reference.}
The $m=1$ specialization is the classical leading equivalent
$a(n)\asym\sqrt{\pi n}$. It does not establish the stronger $O(1)$
error term in A073047.

\clearpage
